\PassOptionsToPackage{dvispnames, table}{xcolor}
\documentclass[aspectratio=169, hyperref={colorlinks=true}]{beamer}
\usepackage[beamer]{prettytex/base}
\usepackage{prettytex/math}
\usepackage{prettytex/gfx}
\usepackage[backend=biber, citestyle=numeric, bibstyle=numeric, hyperref=true, maxbibnames=99]{biblatex}
\usepackage{csquotes}
\usepackage{fontawesome}
\addbibresource{../literature/sources.bib}
\usetheme{oxonian}

% \definecolor{applegreen}{rgb}{0.55,0.71,0.0}
% \definecolor{americanrose}{rgb}{1.0,0.01,0.24}
% \definecolor{aqua}{rgb}{0.0,1.0,1.0}

\title{General Kernel Spectral Methods for Equilibrium Measures \\ \normalsize An MMSC Dissertation}
\titlegraphic{\includegraphics[width=2cm]{theme/logos/ox-circle.eps}}
\author{Peter Julius Waldert}
\institute{Mathematical Institute \\ University of Oxford}
\date{\today}

\begin{document}
  {\setbeamertemplate{footline}{}\frame{\titlepage}}

  \section{Intro}
   {
    \setbeamertemplate{footline}{}
    \begin{frame}{Project Motivation and Goal}
      \vspace{0.4cm}
      \begin{itemize}
        \item Simulate Many-Particle-Systems and find their Equilibrium Distribution.
        \item Interactions through (power-law) Attraction-Repulsion Potentials\footnote{If the repulsive part is larger (so $\beta > \alpha$), there is no equilibium distribution as particles simply continue repelling one another out to infinity.}
              $$K(r) = \frac{r^\alpha}{\alpha} - \frac{r^\beta}{\beta}\quad \text{ with parameters } \quad \alpha, \beta \in \R \backslash \{0\}\,.$$
        \item Each particle $i=1, ..., N$ at position $\vec{x_i} \in \R^d$ and time $t \in \R^+$ follows
              $$\frac{\dd^2 \vec{x_i}}{\ddt^2} = f\left(\norm{\frac{\dd \vec{x_i}}{\ddt}}_2\right) \frac{\dd\vec{x_i}}{\ddt} - \frac{1}{N} \sum_{j=1, i\neq j}^{N} \nabla K\left(\norm{\vec{x_i} - \vec{x_j}}_2\right)\,,$$
              for reference see, for example, \parencite{2020-power-law-kernels, 2021-arbitrary-dimensions}.
              For now, we only consider the case without an external potential $V(\vec{x})$.
      \end{itemize}
    \end{frame}
   }

  \begin{frame}{An Integral Equation}
    The total energy is $$E = \sum_{i=1}^{N} \sum_{j=1, j \neq i}^{N} K\left(\norm{\vec{x_i} - \vec{x_j}}_2\right)\,,$$
    which, in the continuous limit ($N \rightarrow \infty$), becomes
    $$E = \frac{1}{2} \iint K\left(\norm{\vec{x} - \vec{y}}_2\right) \,\dd\rho(\vec{x})\,\dd\rho(\vec{y})\,,$$ where $\dd\rho = \rho(\vec{x})\dd\vec{x}$ is a measure chosen such that $$M = \int \dd\rho = \int_{\supp(\rho)} \rho(\vec{x}) \,\dd\vec{x} = 1\,.$$
  \end{frame}

  \begin{frame}{Simulation Results}
    \begin{figure}
      \centering
      \includegraphics[width=0.64\linewidth]{figures/screenshot5.png}
      \caption{The positional distribution approached by $N = 250$ particles.}
    \end{figure}
  \end{frame}

  \begin{frame}{Spectral Methods}
    \begin{itemize}
      \item In order to find $\rho(x)$, we consider the following ansatz
            $$\rho(\vec{x}) = \sum_{n=0}^{\infty} f_n p_n(\vec{x})\,,$$
            with which we construct a spectral method for the numerical solution of the above integral equation.
            % \item Using the Jacobi polynomial basis (see next slide).
      \item Minimization routine of $E$ over coefficients in $\rho$, as a subroutine of outer minimisation over the bounds of the box (simpler case: use $[-r, r]$, $r \in \R^+$).
      \item By construction, we find that we do not need the inner optimisation routine.
      \item The outer minimisation can be performed using known methods from continuous optimisation.
    \end{itemize}
  \end{frame}

  \begin{frame}{Complete Polynomial Basis}
    Jacobi polynomials $P_n^{(a,b)}(x)$ are orthogonal on $[-1,1]$ w.r.t. the weight function
    \begin{equation*}
      w^{(a,b)}(x)=(1-x)^a (1+x)^b,
    \end{equation*}
    so they satisfy
    \begin{align*}\label{eq:orthogonalityconditionJacobi}
      \int_{-1}^1(1-x)^a(1+x)^bP_n^{(a,b)}P_m^{(a,b)}\ddx = \frac{2^{a+b+1} \Gamma (a+n+1) \Gamma (b+n+1)}{n! (a+b+2 n+1) \Gamma (a+b+n+1)} \delta _{n,m},
    \end{align*}
    with $a	,b>-1$, which uniquely determines $P_n^{(a,b)}(x)$. The special case of $a=b$ corresponds to the ultraspherical or Gegenbauer polynomials, while the case $a=b=0$ corresponds to the Legendre polynomials \cite[Table 18.3.1]{2018-nist}.

    \begin{itemize}
      \item This basis yields a \textbf{sparse}, and in particular, \textbf{banded} operator.
    \end{itemize}
  \end{frame}

  \begin{frame}{Results from the Spectral Solver}
    \vspace{0.4cm}
    Using the spectral method, we try to immediately solve for the distribution function in the continuous limit (here, $d = 1$ dimension).
    \vspace{0.4cm}

    \begin{columns}
      \begin{column}{0.5\textwidth}
        \begin{figure}[H]
          \centering
          \includegraphics[width=0.9\linewidth]{figures/a2b1-5-measure.pdf}
          \caption*{Equilibrium distribution for $\alpha = 2$, $\beta = 1.5$}
        \end{figure}
      \end{column}
      \begin{column}{0.5\textwidth}
        \begin{figure}[H]
          \centering
          \includegraphics[width=0.9\linewidth]{figures/a35b16_supplement_3k.pdf}
          \caption*{Equilibrium distribution for $\alpha = 3.5$, $\beta = 1.6$}
        \end{figure}
      \end{column}
    \end{columns}
  \end{frame}

  \begin{frame}{Two Dimensional Results}
    \vspace{0.4cm}
    The same is possible for $d = 2$ or more dimensions.
    \begin{figure}[H]
      \centering
      \includegraphics[width=0.5\linewidth]{figures/spectral-solution-2d.pdf}
      \caption*{The radially symmetric equilibrium distribution $\rho(\vec{x})$ for $\alpha = 1.2$, $\beta = 0.1993$ in $d = 2$ dimensions.}
    \end{figure}
  \end{frame}

  \begin{frame}{Upcoming Work}
    \begin{itemize}
      \item Establish theory behind spectral methods in the Jacobi basis.
      \item Use kernel expansions to construct an equilibrium measure method on $[-1, 1]$.
      \item Numerically solve for $\rho(x)$ using a fully implemented spectral solver\footnote{This is the first numerical method for particle distribution problems!}.
      \item Implement a particle simulator for the same kernel $K$ (mostly done).
      \item Compare results of the two methods.
    \end{itemize}
    Further goals include
    \begin{itemize}
      \item Considering the $d$-dimensional unit ball domain ($d > 1$).
      \item Optionally: Extend to l-Morse potentials.
    \end{itemize}
  \end{frame}

  \begin{frame}{}
    Questions?
  \end{frame}

  \begin{frame}[allowframebreaks]
    \frametitle{Bibliography}
    \printbibliography[heading=bibnumbered]
  \end{frame}
\end{document}
